% RQ2 Abstract: Governance Capture Mitigation

Governance capture---where a small group gains disproportionate control over collective decision-making---represents an existential threat to decentralized organizations. Token-weighted voting, the dominant mechanism in DAOs, is inherently plutocratic, enabling wealthy actors to dominate governance outcomes.

We present a systematic comparison of three mitigation strategies using multi-agent simulation: (1) vote caps that limit maximum individual influence, (2) quadratic voting that applies diminishing returns to token power, and (3) velocity penalties that restrict rapid token accumulation for governance purposes.

Through \experimentruns{} simulation runs across 27 parameter configurations, we evaluate each strategy's effectiveness at reducing whale influence while maintaining governance throughput. Our results reveal fundamental tradeoffs: vote caps effectively limit peak influence but create cliff effects; quadratic mechanisms smooth the power curve but remain susceptible to sybil attacks; velocity penalties address rapid accumulation but add implementation complexity.

We identify optimal parameter ranges for each mechanism and discuss hybrid approaches combining multiple strategies. Our findings provide evidence-based guidance for DAO designers seeking to balance decentralization with operational efficiency.
